% GENERATED FILE. Edit canonical lessons, then run npm run generate:books.
% canonical-source-hash: fnv1a64:b816c1be57bfb95e
% canonical-lessons: MW-C152-vo, MW-C152-yo-vo, MW-R152-he-and-that

\chapter{He and That}
\label{ch:mw-a2-he-and-that}
\hlchaptermodality{152}

\begin{chapteropening}
\textbf{By the end of this chapter:} \emph{I can say who he is, point at this and that, and still contrast, offer a choice and give a reason.}

Complete MW-R152-he-and-that: say who he is, tell this from that, and contrast, choose and give a reason with them.
\end{chapteropening}

\section[\texorpdfstring{\mw{वो} (vo)}{वो (vo)}]{\mw{वो} — he, the man you talk about}

\label{lesson:MW-C152-vo}



\begin{quote}
Say how you are: \textbf{\mw{हूं} \mw{ठीक} \mw{हूं।}} (\emph{hū̃ ṭhīk hū̃.}) Then ask the polite question: \textbf{\mw{आप} \mw{कैसो} \mw{हो}?} (\emph{āp kaiso ho?})

So far you talk about yourself, \textbf{\mw{हूं}}, and to the person in front of you, \textbf{\mw{आप}}. Now a third person: someone you talk \emph{about}.
\end{quote}

\begin{grammarlens}[title={Grammar Lens: \mw{वो}}]
\begin{quote}
\textbf{\mw{वो} \mw{राम} \mw{है।}}
\end{quote}

\begin{quote}
\emph{vo rām hai.} — He is Ram.
\end{quote}

\textbf{\mw{वो}} (\emph{vo}) is \emph{he}: a man you are talking about, not to. He comes first, and \textbf{\mw{है}} ends the sentence, as it does in \textbf{\mw{म्हारो} \mw{नाम} \mw{राम} \mw{है।}}

\begin{quote}
\textbf{\mw{वो} \mw{बीमार} \mw{है।}}
\end{quote}

\begin{quote}
\emph{vo bīmār hai.} — He is ill.
\end{quote}

Old Marwari already had \emph{vo}, \textquotedblleft{}that\textquotedblright{}: CDIAL lists it under Sanskrit \emph{asau}, \textquotedblleft{}that\textquotedblright{} (CDIAL 972). This book keeps \textbf{\mw{वो}} for a man, and names a woman by a word you know: \textbf{\mw{मां}}, \textbf{\mw{बहन}}, \textbf{\mw{छोरी}}.
\end{grammarlens}

\subsection*{Your turn}
Say these aloud:

\begin{itemize}
\raggedright
  \item \textbf{\mw{वो} \mw{छोरो} \mw{है।}} — \emph{vo chhoro hai.} — He is a boy.
  \item \textbf{\mw{वो} \mw{बीमार} \mw{है} \mw{पण} \mw{हूं} \mw{ठीक} \mw{हूं।}} — \emph{vo bīmār hai paṇ hū̃ ṭhīk hū̃.} — He is ill, but I am fine.
  \item \textbf{\mw{हूं} \mw{राजी} \mw{हूं} \mw{क्यूंकै} \mw{वो} \mw{अठै} \mw{है।}} — \emph{hū̃ rājī hū̃ kyū̃kai vo aṭhai hai.} — I am happy because he is here.
\end{itemize}

\begin{tcolorbox}[breakable,colback=teal!4,colframe=teal!35!black,title={Before you move on}]
Who is \textbf{\mw{वो}}? (A man you are talking about.) Which word ends \textbf{\mw{वो} \mw{बीमार} \mw{है।}}? (\textbf{\mw{है}}.)
\end{tcolorbox}
\section[\texorpdfstring{\mw{यो} · \mw{वो} (yo · vo)}{यो · वो (yo · vo)}]{\mw{यो} and \mw{वो} — this one here, that one there}

\label{lesson:MW-C152-yo-vo}



\begin{quote}
Say \emph{he is a boy}. Then the near word you know, \textbf{\mw{यो}} (\emph{yo}), \textquotedblleft{}this\textquotedblright{}. And offer the old choice once more: \textbf{\mw{चाय} \mw{अर} \mw{पाणी}?} (\emph{chāy ar pāṇī?})
\end{quote}

\begin{grammarlens}[title={Grammar Lens: \mw{वो} points away}]
\textbf{\mw{वो}} also points at a thing away from you: \emph{that}. \textbf{\mw{यो}} is the near one, \emph{this}.

\begin{quote}
\textbf{\mw{यो} \mw{म्हारो} \mw{घर} \mw{है।}}
\end{quote}

\begin{quote}
\emph{yo mhāro ghar hai.} — This is my house.
\end{quote}

\begin{quote}
\textbf{\mw{वो} \mw{दुकान} \mw{है।}}
\end{quote}

\begin{quote}
\emph{vo dukān hai.} — That is a shop.
\end{quote}

So \textbf{\mw{वो}} is \emph{he} or \emph{that}: a man, or a thing over there.
\end{grammarlens}

\subsection*{Your turn}
Offer two things the way you offered tea or water: \textbf{\mw{यो} \mw{अर} \mw{वो}?} (\emph{yo ar vo?}), this or that?

Say these aloud:

\begin{itemize}
\raggedright
  \item \textbf{\mw{यो} \mw{पाणी} \mw{है।}} — \emph{yo pāṇī hai.} — This is water.
  \item \textbf{\mw{वो} \mw{चाय} \mw{है।}} — \emph{vo chāy hai.} — That is tea.
  \item \textbf{\mw{यो} \mw{अर} \mw{वो}?} — \emph{yo ar vo?} — This or that?
  \item \textbf{\mw{वो।}} — \emph{vo.} — That one.
\end{itemize}

\begin{tcolorbox}[breakable,colback=teal!4,colframe=teal!35!black,title={Before you move on}]
Which is near, \textbf{\mw{यो}} or \textbf{\mw{वो}}? (\textbf{\mw{यो}}.) What else can \textbf{\mw{वो}} be? (\emph{He}.)
\end{tcolorbox}
\section[(dialogue)]{He and that — a review}

\label{lesson:MW-R152-he-and-that}



\begin{quote}
Say \emph{he is ill}.
\end{quote}

\subsection*{Your turn}
Say these aloud:

\begin{itemize}
\raggedright
  \item \emph{he is Ram}; \emph{he is a boy}; \emph{he is here}
  \item \emph{this is my house}; \emph{that is a shop}; \emph{this or that?}
  \item \emph{he is ill, but I am fine}
  \item \emph{I am happy because he is here}
  \item \emph{tea or water?}, and answer with one of them
\end{itemize}

\begin{tcolorbox}[breakable,colback=teal!4,colframe=teal!35!black,title={Before you move on}]
\emph{He is here}? (\textbf{\mw{वो} \mw{अठै} \mw{है।}}) \emph{That is a shop}? (\textbf{\mw{वो} \mw{दुकान} \mw{है।}}) \emph{Tea or water?} (\textbf{\mw{चाय} \mw{अर} \mw{पाणी}?}) \emph{He is ill, but I am fine}? (\textbf{\mw{वो} \mw{बीमार} \mw{है} \mw{पण} \mw{हूं} \mw{ठीक} \mw{हूं।}})
\end{tcolorbox}
